\chapter{Metodología}
La metodología desarrollada en este trabajo de investigación se estructuró sistemáticamente en tres fases secuenciales para garantizar su rigor y reproducibilidad técnica como se ilustra en la Figura \ref{fig:metodologia_flowchart}.

En primer lugar, se llevó a cabo la adquisición de datos, la cual consistió en la inoculación controlada de los frutos y la captura temporal de sus firmas hiperespectrales. La segunda fase se centró en el preprocesamiento computacional de las imágenes, abarcando la corrección radiométrica, la segmentación espacial de las regiones de interés (ROI) y la posterior extracción de métricas estadísticas representativas (reflectancia media y desviación estándar) para cada muestra. Finalmente, la tercera fase consistió en el diseño e implementación de una estrategia de optimización basada en algoritmos metaheurísticos, cuya aptitud fue evaluada mediante un modelo de Máquina de Vectores de Soporte (SVM). El objetivo de esta última etapa fue reducir drásticamente la dimensionalidad del espacio de datos e identificar el subconjunto óptimo de bandas espectrales para la detección de \textit{Aspergillus flavus}.
\begin{figure}[H]
    \centering
    \includegraphics[width=0.6\textwidth]{metodology_flow_chart.jpg}
    \caption{Diagrama de flujo de la metodología computacional propuesta}
    \label{fig:metodologia_flowchart}
\end{figure}

\section{Adquisición de Datos e Inoculación}
\subsection{Material Biológico y Procedencia}
Para el desarrollo de esta investigación, se utilizaron muestras de higos frescos (Ficus carica L.) correspondientes a la variedad Calabacita. Con el fin de garantizar la uniformidad fisiológica del lote de estudio, todas las muestras fueron recolectadas simultáneamente en la "Finca La Orden-Valdesequera", una instalación agro-tecnológica operada por el Centro de Investigaciones Científicas y Tecnológicas de Extremadura (CICYTEX), ubicada en la región de Guadajira, Badajoz, España.

\subsection{Diseño Experimental y Proceso de Inoculación}
El protocolo de inoculación fue ejecutado bajo condiciones controladas de laboratorio por ingenieros agrícolas del CICYTEX. Para el diseño del experimento, la población total de frutos se distribuyó en cuatro bandejas independientes, albergando cada una entre 36 y 40 higos. Una de estas bandejas se mantuvo libre de patógenos para fungir como el grupo de control sano (Clase 0).

Las muestras de las tres bandejas restantes fueron sometidas a un proceso de inoculación artificial, el cual consistió en sumergir la porción inferior de cada fruto durante un lapso de 4 a 5 segundos en una solución fúngica. Se utilizó específicamente la cepa Aspergillus flavus HG 144M, provista por la colección del grupo de investigación CAMIALI de la Universidad de Extremadura.
Para evaluar la progresión de la enfermedad, las soluciones se prepararon en tres concentraciones distintas: $10^3 UFC/ml$, $10^5 UFC/ml$ y $10^7 UFC/ml$. Sin embargo, para los alcances del presente trabajo computacional, centrado en la clasificación binaria, el conjunto de datos de trabajo se construyó utilizando únicamente el grupo de control sano (Clase 0) y el grupo inoculado con la infección más severa de $10^7 UFC/ml$ (Clase 1).

\subsection{Condiciones de Incubación y Monitoreo Temporal}
Entre las sesiones de adquisición de imágenes, las muestras reposaron en una cámara de incubación bajo un ambiente estrictamente controlado. La cámara se configuró para mantener una temperatura constante de 25°C y una humedad relativa de entre 80\% y 90\%, proporcionando las condiciones ambientales óptimas para la proliferación fúngica.

El monitoreo de las muestras se realizó mediante la captura de imágenes hiperespectrales cada 24 horas durante un periodo de cinco días consecutivos posteriores a la inoculación. La adquisición se efectuó empleando una cámara hiperespectral line-scan (pushbroom) modelo SPECIM FX10. El sensor de este dispositivo captura la reflectancia en el rango espectral visible e infrarrojo cercano (VNIR, 400 nm a 1000 nm), ofreciendo una resolución espectral de 448 bandas y una resolución espacial de 1024 píxeles de ancho por cada línea de escaneo \cite{cruz-carrasco_detection_2024}.

\section{Preprocesamiento de Imágenes}
El procesamiento de los datos espaciales y espectrales, así como la extracción inicial de características, se llevó a cabo apoyándose en la biblioteca Spectral Python (SPy) para la manipulación y análisis algorítmico de los cubos hiperespectrales. Como primer paso en el flujo de trabajo, se procedió a la lectura e interpretación de los archivos de encabezado (.hdr) y los archivos de datos binarios (.raw) generados por el sensor SPECIM FX10, lo que permitió reconstruir digitalmente las escenas adquiridas.
Como primer paso en el flujo de trabajo, se procedió a la lectura e interpretación de los archivos de encabezado (.hdr) y los archivos de datos binarios (.raw) generados por el sensor SPECIM FX10, lo que permitió reconstruir digitalmente las escenas adquiridas.

Dado que los valores capturados inicialmente por la cámara representan intensidades lumínicas brutas, las cuales están inherentemente sujetas a fluctuaciones de la fuente de iluminación y al ruido térmico del sensor, fue indispensable aplicar una corrección radiométrica. Los valores de intensidad bruta se convirtieron a reflectancia relativa aplicando referencias claras y oscuras adquiridas junto con cada escena, de acuerdo con la siguiente ecuación\cite{morales_laboratory_2022}:
\begin{equation}
    R=\dfrac{I_{raw}-I_{dark}}{I_{white}-I_{dark}}
\end{equation}

Donde $I_{raw}$ corresponde a la intensidad de la imagen bruta, $I_{dark}$ representa la referencia oscura (capturada con el obturador cerrado para medir la corriente oscura del sensor), e $I_{white}$ es la referencia blanca (obtenida mediante un panel de calibración de reflectancia máxima).

Posterior a la calibración, y con el fin de procesar la información espectral de cada fruto de manera individualizada y libre del ruido de fondo, se procesaron un total de 18 escenas hiperespectrales completas. El aislamiento espacial de las Regiones de Interés (ROI) se efectuó a través de una segmentación manual utilizando el software de anotación de código abierto LabelMe (v5.8.3)\cite{wada_wkentarolabelme_2021}.
Este meticuloso proceso dio como resultado la generación de 560 máscaras de segmentación individuales. Consecuentemente, el conjunto de datos de trabajo final quedó constituido por estas 560 muestras independientes, distribuidas en 277 muestras pertenecientes a la clase de control sana (Clase 0) y 283 muestras correspondientes a la clase inoculada con la concentración severa (Clase 1).

\begin{figure}[H]
    \centering
    \subfloat[]{\includegraphics[width=0.4\textwidth]{sample_figs_tray.png}}
    \hspace{1mm}
    \subfloat[]{\includegraphics[width=0.4\textwidth]{sample_figs_masks.png}}
    \caption{Charola de higos antes de la segmentación (a) y mascaras de higos después de la segmentación (b)}
    \label{fig:segmentacion}
\end{figure}

\section{Extracción de Características y Construcción del Vector}
La información contenida en las imágenes hiperespectrales es masiva y altamente redundante entre bandas contiguas. Para modelar el problema de clasificación, fue necesario transformar los datos puros de los píxeles dentro de cada Región de Interés (ROI) en descriptores matemáticos \cite{sawant_survey_2020}.
Esta característica conlleva a la llamada \textit{maldición de la dimensionalidad} o fenomeno de Hughes, donde la presición de la clasificación tienede a decrecer conforme el número de bandas espectrales se incrementa dado un número limitado de muestras de entrenamiento \cite{hughes_mean_1968}.

Con el objetivo de capturar tanto el impacto fisiológico interno (cambios en la pigmentación y contenido de agua) como las alteraciones en la textura superficial del fruto, se diseñó una estrategia de construcción de características basada en dos momentos estadísticos:

\begin{enumerate}
    \item Reflectancia Media ($\mu$): Caracteriza la tendencia central de la firma espectral, indicando la cantidad de luz promedio que el tejido absorbe o refleja en una longitud de onda específica \cite{nagasubramanian_hyperspectral_2018}.
    \item Desviación Estándar ($\sigma$): Cuantifica la variabilidad espacial o heterogeneidad en la superficie del fruto. Esta métrica es de suma importancia biológica, ya que la infección por Aspergillus flavus no es uniforme; la formación de micelio y la esporulación alteran la rugosidad superficial del higo, generando patrones de dispersión de luz irregulares que son capturados por esta medida de dispersión \cite{hossain_color_2019}.
\end{enumerate}

Formalmente, para una imagen hiperespectral con n bandas (donde $n=448$), el conjunto completo de características extraídas para la $i$-ésima muestra se definió mediante el siguiente vector concatenado:
\begin{equation}
    X_i=[\mu_i,1,\ldots \mu_i,n,\sigma_i,1,\ldots,\sigma_i,n]
\end{equation}

Este enfoque generó un vector inicial de 896 dimensiones por muestra (448 medias y 448 desviaciones estándar). Para mitigar la maldición de la dimensionalidad, el espacio de búsqueda de optimización se estructuró como un problema de selección de bandas.
Bajo este esquema, la selección de una banda específica j implica la inclusión simultánea e indisoluble de su media ($\mu_j$) y su desviación estándar ($\sigma_j$).
Por lo tanto, el algoritmo metaheurístico opera sobre un espacio de búsqueda binario de n dimensiones (donde 1 indica selección y 0 exclusión), pero la dimensionalidad real del vector de características que evalúa el clasificador es de $ 2\times{k} $, siendo $k$ el número objetivo de bandas seleccionadas ($k\in\{3, 5, 6\}$).

El problema de optimización se definió formalmente como la búsqueda del vector binario $B=\{b1,...,bn\}$ que maximice la función de aptitud, se estableció el F1-Score como la función objetivo a maximizar:

\begin{equation}
    F1=2 \times{\frac{Precision\times{Recall}}{Precision+Recall}}
\end{equation}

Se eligio el F1-Score como la función de aptitud debido a su habilidad para balancear presición y sensibilidad. Es critico penalizar falsos negativos (higos infectados clasificados erroneamente como sanos), ya que representan un alto riesgo para la salud. Ademas, el F1-Score provee una métrica robusta ya que mitiga efectivamente el sesgo potencial introducido por la pequeña diferencia entre la cantidad de higos en Clase 0 y la Clase 1 en el conjunto de datos.

\section{Proceso de Optimización (Diseño Metaheurístico)}
\begin{itemize}
    \item Se implementaron GA y PSO para buscar un vector binario que maximice el rendimiento de clasificación, con tamaños de subconjuntos objetivo de 3, 5 y 6 bandas \cite{eiben_introduction_2015}.
    \item Para mitigar la carga computacional iterativa, se diseñó una estrategia de evaluación en dos etapas: 
    \begin{itemize}
        \item \textbf{Evaluador rápido:} Un modelo SVM ligero configurado para estimar la aptitud de los individuos rápidamente y guiar la búsqueda del algoritmo \cite{nagasubramanian_hyperspectral_2018}.
        \item \textbf{Evaluador exhaustivo:} Un ajuste fino de hiperparámetros utilizado únicamente sobre el mejor subconjunto encontrado al final de la optimización \cite{nagasubramanian_hyperspectral_2018}. Antes del entrenamiento, los datos fueron normalizados con StandardScaler.
    \end{itemize}
\end{itemize}
